\section{Discussion and conclusions}\label{sec:discussion}

\subsection{What the index indicates}
An index is only useful if a high value means something one can act on. A high \Fd\ marks a crop and place where
three things coincide:
\begin{itemize}
\item the growing area is already under heat or water stress;
\item the crop loses yield quickly under that stress;
\item people depend on the crop for nutrients.
\end{itemize}
Read this way, the results say five things.
\begin{enumerate}
\item \textbf{The most vulnerable crop is not the most heat-sensitive one.} Soybean has the steepest literature heat
slope, yet ranks low: its monthly flowering temperature stays below its threshold, and it is eaten mainly as oil.
Wheat is only moderately sensitive, but it is grown in the hottest and driest plains of the panel and is the
largest single food source of zinc and iron in most panel countries. Vulnerability is about \emph{overlap}, which
neither a yield study nor a nutrient study shows on its own.
\item \textbf{The nutrient at stake is iron.} Iron is the only nutrient whose national supply falls below the main
requirement in any panel country (\nIronShort\ of \nAdqCountries). The crops that supply most iron are among those
whose climate stress is highest or rising fastest: wheat everywhere, and pulses in South Asia and East Africa.
Climate risk therefore falls on a nutrient that is already short, not on calories. This rests on the requirement
chosen (Section~\ref{sec:adequacy}).
\item \textbf{Concentrated diets amplify the signal.} Where a few crops supply most nutrients, a loss in one
vulnerable crop removes a larger share of a nutrient. India, for example, gets \divIndIronTopShare\% of its iron from
wheat, and its energy supply is spread over an effective \divIndKcal\ food items against \divWorldKcal\ for the
world as a whole.
\item \textbf{Irrigation is the hidden buffer.} Rice scores low largely because most of its area is irrigated,
which the water measure counts as unstressed. Groundwater depletion or reservoir shortfalls would move rice up the
ranking, and monthly data cannot see the hot days that matter for rice.
\item \textbf{The future changes magnitudes more than the order.} Under one model and one scenario, the crop
ranking in 2041--2060 is largely the same, but vulnerability rises for wheat and common bean. The increase comes
mostly from heat, while today's scores come mostly from water. Pulses, a key protein and iron source for
vegetarian diets, rise fastest relative to their present level (\grpPulsePct\%).
\end{enumerate}

\subsection{What it points to}
As a screening tool, the index points to where adaptation would protect most nutrition:
\begin{itemize}
\item \textbf{wheat in India:} the present score is mainly water, but the projected increase comes almost entirely
from heat during flowering, so heat tolerance and sowing dates matter as well as water;
\item \textbf{wheat in Russia:} both today's score and most of the projected increase are water-driven;
\item \textbf{wheat in the USA:} the increase is shared between heat and water;
\item \textbf{rainfed maize in southern Africa:} drought management;
\item \textbf{beans in tropical Brazil and East Africa:} heat-tolerant varieties;
\item \textbf{on the demand side:} diversifying iron sources, so that one crop's bad year does not decide iron
supply.
\end{itemize}
These are directions suggested by where scores are high, not tested interventions.

\subsection{What the four extensions add}
\begin{itemize}
\item \textbf{Parameter evidence (Section~\ref{sec:parevidence}).} The headline survives the audit but its scope
narrows. Wheat is first in the full index and also first among the seven crops whose parameters are supported, but
rice and rapeseed cannot be placed on supported heat parameters, and bean in India and Myanmar does not match its
parameter source. The full nine-crop ranking is therefore an \emph{assumption-dependent} result.
\item \textbf{Future literature (Section~\ref{sec:futurelit}).} Direction and mechanism agree with published studies
where they are most comparable (wheat in India, beans in Tanzania, mustard in India) and differ where studies include
CO$_2$ fertilisation or adaptation, or where our index cannot see the mechanism (rice, irrigated crops, season
shortening). The studies include evidence that challenges rising vulnerability, and the one climate model we use is
described by one study as the least pessimistic.
\item \textbf{CO$_2$ and nutrient quality (Section~\ref{sec:co2}).} The concentration channel, as an experimental
stress test with partial coverage, lowers national zinc by more than the climate yield channel in all
\cotwoNCountries\ panel countries and iron in \cotwoNQualGtQtyIron, and does so even where the climate scenario removes no
modelled supply. Calories are not estimated. Nothing here is a forecast.
\item \textbf{Diversity after losses (Section~\ref{sec:divloss}).} Losing part of a dominant staple lowers total
supply and can raise the effective number of food sources. Diversity must be read with total supply and adequacy,
never alone.
\end{itemize}

\subsection{Supported conclusions and their strength}
\begin{table}[H]\centering\small
\caption{Conclusions and how strongly the evidence supports them.}\label{tab:conclusions}
\begin{tabular}{p{6.3cm}p{2.0cm}p{6.4cm}}
\toprule
Conclusion & Support & Evidence and caveat \\
\midrule
Wheat is the most vulnerable of the nine crops & strong, conditional & first in \robTopFfour/\robNFfour\ F4 robustness rankings and \mcTopShare\% of random weightings; exceptions only under empirical-only sensitivity. Conditional on proxies: first among the seven crops with supported parameters (rice and rapeseed unavailable) \\
Order of positions 2--9 & weak & moves with weighting and pulse parameters \\
Rice and rapeseed positions & not established & heat slopes unsupported; present scores are almost entirely water \\
Direction of future change agrees with published studies where comparable & moderate, partial & wheat in India, bean in Tanzania, mustard in India; differs where CO$_2$ or adaptation is included; \litNPairsNoStudy\ of \litNPairs\ pairs have no study \\
Elevated-CO$_2$ quality channel lowers zinc, iron and protein & stress test only & covered crops only, no energy or fat, experiment-based step, hybrid reference; not a forecast \\
Crop loss can raise food-source diversity while lowering supply & strong (arithmetic) & predicts the sign in every single-item case; read with adequacy \\
Water dominates today's scores, heat the projected increase & moderate & \presentWaterShare\% water today, \futHeatShare\% heat in the change; water ratio has no soil storage; one climate model \\
Iron is the nutrient closest to shortfall & conditional & \nIronShort/\nAdqCountries\ countries under the main requirement, \nIronShortAlt\ under the alternative \\
Index associated with yield-shock frequency & weak & positive in 59 pairs, null in the 26-pair mid-evaluation sample; association only \\
Future yield and nutrient losses & screening only & literature and empirical responses disagree for about half the pairs \\
Hyderabad's risk is more structural than climatic & moderate & district anaemia with adequate energy; rice irrigated; based on state-urban and district data, not city intake \\
\bottomrule
\end{tabular}
\end{table}

\subsection{Conclusion}
Across nine crops, the Climate--Nutrition Vulnerability Index identifies \MakeLowercase{\topEq} as the crop whose
climate stress most overlaps with its nutritional role. This holds across \robTopFfour\ of \robNFfour\ robustness
rankings and is stable to 2041--2060 under one climate model. Cereals as a group are most vulnerable, pulses rise
fastest, and the nutrient most at stake is iron. For India the index points to wheat, chickpea and common bean; for
Hyderabad it points to a rice-based diet whose main problem, iron density, exists independently of climate. The
index is most useful as a map of where climate stress and dietary dependence coincide, which is where adaptation
and diet diversification would protect nutrition most. It is not a forecast of yields or diets.

The audit of parameters, the comparison with future-impact studies, the CO$_2$ stress test and the diversity analysis
sharpen this reading without overturning it. Wheat's first place does not depend on the unsupported parameters, but the
positions of rice and rapeseed cannot be established with them, several studies that include CO$_2$ fertilisation or
adaptation project gains where we project losses, quality changes in grain compound the quantity loss for zinc, iron
and protein, and a loss in a dominant staple can make food sources look more even while supply falls. None of this
means the model is the most accurate available; the unresolved evidence gaps are listed in
Sections~\ref{sec:parevidence}--\ref{sec:co2} and the human review tasks in Section~\ref{sec:repro}.
